\documentclass[11pt]{article}
\input{style-sujet}
\usepackage{array}

\begin{document}

\entetesujet{Sujet bac 2012 -- Série D}

% ====================== EXERCICE 1 ======================
\exo{1}{4 points}

On considère la série statistique à double variable $X$ et $Y$ définie par le tableau ci-après :

\begin{center}
\renewcommand{\arraystretch}{1.3}
\begin{tabular}{|c|c|c|c|c|c|}
\hline
$X$ & $-2$ & $0$ & $1$ & $a$ & $4$ \\
\hline
$Y$ & $-10$ & $-8$ & $b$ & $0$ & $12$ \\
\hline
\end{tabular}
\end{center}

\begin{enumerate}
  \item Déterminer les réels $a$ et $b$ pour que le point moyen $G$ du nuage statistique, ait pour coordonnées $(1\,;-2)$.
  \item Dans la suite, on prendra $a = 2$ et $b = -4$.
  \begin{enumerate}[label=\alph*.]
    \item Représenter graphiquement les points du nuage de cette série statistique.
    \item Déterminer l'équation de la droite de régression de $X$ en $Y$.
    \item Calculer le coefficient de corrélation linéaire entre $X$ et $Y$, puis interpréter le résultat.
  \end{enumerate}
\end{enumerate}

% ====================== EXERCICE 2 ======================
\exo{2}{4 points}

\begin{enumerate}
  \item Résoudre dans l'ensemble $\Cb$ des nombres complexes, l'équation :
  \[ (E) : z^2 + 8\sqrt{3} - 8i = 0 \]
  \begin{enumerate}[label=\alph*.]
    \item En utilisant la forme trigonométrique.
    \item En utilisant la forme algébrique. On pourra admettre que $8 + 4\sqrt{3} = \left(\sqrt{2} + \sqrt{6}\right)^2$.
  \end{enumerate}
  \item Placer les images des solutions $z_1$ et $z_2$ de $(E)$ sur un cercle trigonométrique.
  \item Déduire de ce qui précède, la valeur exacte de $\cos\!\left(\dfrac{5\pi}{12}\right)$ et $\sin\!\left(\dfrac{5\pi}{12}\right)$.
\end{enumerate}

% ====================== PROBLÈME ======================
\prob{12 points}

\partie{A}
\begin{enumerate}
  \item Résoudre l'équation différentielle : $y'' + 2y' + y = 0$.
  \item Déterminer la solution particulière $u$, sachant que $u(0) = 1$ et $u'(0) = 0$.
\end{enumerate}

\partie{B}
Soit $f$, la fonction numérique de la variable réelle $x$ définie par :
\[ f(x) = \begin{cases} (x + 1)e^{-x} & \text{si } x \le 0 \\[0.5ex] 1 - 2x + x\ln x & \text{si } x > 0 \end{cases} \]
On désigne par $(C)$ la courbe représentative de $f$ dans le repère orthonormé $(O,\vi,\vj)$ du plan d'unité graphique : 2 cm.

\begin{enumerate}[start=3]
  \item Préciser l'ensemble de définition de $f$.
  \item Étudier la continuité et la dérivabilité de $f$ en $x = 0$.
  \item Étudier les variations de $f$. On dressera un tableau de variation de $f$.
  \item Pour $x \le 0$, déterminer les coordonnées du point d'intersection de la courbe $(C)$ avec l'axe des abscisses et écrire une équation cartésienne de la tangente $(T)$ à $(C)$ en ce point.
  \item Démontrer que l'équation $f(x) = 0$ admet une solution unique $\alpha \in ]6\,;7[$. On ne demande pas de calculer $\alpha$.
  \item \begin{enumerate}[label=\alph*.]
    \item Étudier les branches infinies à $(C)$.
    \item Tracer la courbe $(C)$ de $f$ et la droite $(T)$.
  \end{enumerate}
\end{enumerate}

\partie{C}
Soit $h$ la fonction numérique de la variable réelle $x$, définie par :
\[ \forall x \in \R \quad h(x) = -f(x) \]

\begin{enumerate}[start=9]
  \item \begin{enumerate}[label=\alph*.]
    \item Dresser le tableau de variation de $h$.
    \item Tracer la courbe $(C')$ représentative de $h$ dans le même repère que $(C)$ de $f$.
    \item Calculer en cm$^2$, l'aire $\mathcal{A}$ du domaine $(D)$ limité par les courbes $(C)$ ; $(C')$ et les droites d'équations $x = -1$ ; $x = 0$.
  \end{enumerate}
\end{enumerate}

\end{document}
